% Canonical Paper IV section.
Contextual residuals measure semantic distinguishability.  An execution uses
concrete data structures, schedules, memories, and communication links.  The
two layers meet only after a machine model says where a complete residual state
is represented.

\subsection{Logical cuts and live configurations}

Let $\Gamma$ be a typed acyclic computation network.  A logical cut separates
operations already licensed by causality from operations not yet licensed.  A
set of nodes that have ever been computed may record causal progress, but it
does not say which values are currently present.

\begin{definition}[Operational machine model]
An operational model $\cM$ specifies:
\begin{enumerate}
  \item the admissible input family and uniformity convention;
  \item value, instruction, and constant encodings;
  \item atomic compute, erase, move, communicate, checkpoint, and recompute
  actions;
  \item legal live configurations and their sizes;
  \item the cost of each action and the processor or memory topology;
  \item whether input, output, constants, tables, and code count toward space.
\end{enumerate}
\end{definition}

An execution is a legal trace
\begin{equation}
  L_0\xrightarrow{a_1}L_1\xrightarrow{a_2}\cdots
  \xrightarrow{a_T}L_T,
  \label{eq:live-configuration-trace}
\end{equation}
where $L_t$ lists the values and control records currently materialized.
Write $S_\sigma(t)=\operatorname{size}_{\cM}(L_t)$.

\begin{definition}[Resource vector]
For a realization $\Gamma$ and schedule $\sigma$, define
\begin{equation}
 \mathbf C_{\cM}(\Gamma,\sigma)
 =\bigl(
 L_{\mathrm{desc}},W,T,D_{\mathrm{caus}},S_{\max},\ST,Q
 \bigr),
 \label{eq:resource-vector}
\end{equation}
where
\[
 S_{\max}=\max_t S_\sigma(t),
 \qquad
 \ST=\sum_{t=0}^{T}S_\sigma(t).
\]
Here $L_{\mathrm{desc}}$ is static description size, $W$ weighted work,
$T$ scheduled time, $D_{\mathrm{caus}}$ intrinsic causal depth in the declared
parallel model, and $Q$ communication or I/O.
\end{definition}

These coordinates generally have different units.  A theorem comparing them
must state the simulation, charging, and encoding convention; adding them into
one scalar without such a rule has no invariant meaning.

\subsection{Why completed-node sets do not determine space}

\begin{theorem}[Completed-set insufficiency]
\label{thm:completed-set-insufficient}
Suppose a computation has a positive-size intermediate value that may legally
be retained or erased after its last immediate use.  Then the set of nodes ever
computed does not determine the current workspace.
\end{theorem}

\begin{proof}
Execute through the last use of the intermediate.  In one legal execution keep
its value; in another erase it.  The two executions have computed exactly the
same nodes, but their live configurations differ by a positive-size object.
\end{proof}

With recomputation the discrepancy is stronger: a value may disappear and
later be recreated while the ever-computed set remains unchanged.  Therefore a
monotone chain of down-sets cannot serve as a general pebble or checkpoint
state.

\begin{proposition}[Restricted frontier model]
\label{prop:restricted-frontier}
In a no-recomputation node-value model in which each value is retained exactly
until its last successor has executed, a down-closed completed set $I$ does
determine the live frontier
\[
  \partial^+ I
  =\{u\in I:\text{some successor of $u$ lies outside $I$}\}.
\]
If node $u$ occupies $b(u)$ units, then the corresponding workspace is
$\sum_{u\in\partial^+I}b(u)$, up to the separately declared input, output, and
control conventions.
\end{proposition}

The proposition remains useful, but only inside its restricted retention
policy.

\subsection{Embedding the residual lower bound}

A live configuration may carry much more than the minimal residual class.  It
may include redundant values, caches, checksums, code pointers, or a convenient
structured representation.  Nevertheless, when it is the complete exact online
state, it must distinguish residual classes.

\begin{proposition}[Live-state residual bound]
\label{prop:live-residual-bound}
At a cut $C$, suppose the charged live configuration contains a complete exact
state and is encoded either in a fixed-width $b$-bit slot or by a prefix-free
code of maximum length $b$.  Then
\[
  b\ge\left\lceil\log_2|\cR_C|\right\rceil.
\]
\end{proposition}

\begin{proof}
The live-state encoding is an exact online realization, so
Theorem~\ref{thm:minimal-exact-residual} maps its reachable states onto
$\cR_C$.  Apply the counting argument of
Corollary~\ref{cor:residual-bit-bound}.
\end{proof}

The following quantities must remain distinct:

\begin{center}
\begin{tabular}{L{0.25\textwidth}L{0.65\textwidth}}
\toprule
Quantity & Question answered\\
\midrule
$|\cR_C|$ & How many future behaviors remain semantically distinguishable?\\
$\lceil\log_2|\cR_C|\rceil$ & How many bits select one exact residual class in a fixed-width model?\\
$S_\sigma(t)$ & How large is the materialized configuration in this execution?\\
$L_{\mathrm{desc}}$ & How large is the program, circuit, table, or shared representation?\\
\bottomrule
\end{tabular}
\end{center}

A dense tensor with $D$ field entries, a low-rank factorization of the same
tensor, an index into a decision diagram, and the whole decision diagram are
four different representations.  Their sizes are not determined by the
cardinality of one abstract residual space.

\subsection{Uniform families and Pareto resource geometry}

Complexity belongs to input families rather than to one hard-coded finite map.
Let $F=(F_n)_{n\ge1}$ be a family and let
$\operatorname{Adm}^{\mathrm{unif}}_{\cM}(F)$ be the realization families
produced by one declared uniform constructor.  Define the fixed-size achievable
frontier
\begin{equation}
 \mathfrak C_{\cM}^{\mathrm{unif}}(F;n)
 =\Pareto\left\{
 \mathbf C_{\cM}(\Gamma_n,\sigma_n):
 ((\Gamma_j,\sigma_j))_{j\ge1}
 \in\operatorname{Adm}^{\mathrm{unif}}_{\cM}(F)
 \right\}.
 \label{eq:pareto-resource-geometry}
\end{equation}
Here $\Pareto$ retains nondominated vectors under coordinatewise comparison.
If limits are not attained, a closure may be needed.

Equation~\eqref{eq:pareto-resource-geometry} is a structural framework, not a
machine-independent invariant.  Its value changes with the model, units,
precision, uniformity, preprocessing, and allowed representations.  A robust
comparison across models would require explicit efficient simulations.

\subsection{Three distinct order effects}

The examples of this paper separate three phenomena:

\begin{center}
\begin{tabular}{L{0.25\textwidth}L{0.35\textwidth}L{0.28\textwidth}}
\toprule
Order effect & Meaning & Calibration\\
\midrule
semantic noncommutation & Order changes the induced operator or endpoint & affine add--scale histories\\
rewrite-plan variation & Semantics agree but factorization and costs differ & matrix-chain parenthesization\\
schedule or cut sensitivity & One computation exposes different live or residual widths & OBDD order and checkpointing\\
\bottomrule
\end{tabular}
\end{center}

Only the first already has the contact-curvature interpretation developed in
Paper~I~\cite{Yuan2026AEGFoundations}.  The other two should not be called
curvature until a connection, gauge law, or $2$-cell defect has been supplied.
